\subsection{Oriented graphs and quivers}

Let G be a graph, \((S,F,o,t)\) the underlying quiver, and A an orientation of G. Then, \((S,A,o\mid A,t\mid A)\) is a quiver, called the quiver associated with the oriented graph \((G,A)\) .

Conversely, let \(C = (S, F, o, t)\) be a quiver. Set \(\widetilde{F} = F \times \{-1, 1\}\) and let \(\widetilde{o}, \widetilde{t}\) be the maps from \(\widetilde{F}\) to S defined by

\[
\begin{array}{l l} \widetilde {o} (f, 1) = o (f) & \qquad \qquad \widetilde {o} (f, - 1) = t (f) \\ \widetilde {t} (f, 1) = t (f) & \qquad \qquad \widetilde {t} (f, - 1) = o (f), \end{array}
\]

for every \(f \in \mathrm{F}\) . Then, \(\widetilde{\mathrm{C}} = (\mathrm{S}, \widetilde{\mathrm{F}}, \widetilde{o}, \widetilde{t})\) , equipped with the involution \((f, \varepsilon) \mapsto (f, -\varepsilon)\) of \(\widetilde{\mathrm{F}}\) is a graph, called the graph associated with the quiver C. The set \(\mathrm{A} = \mathrm{F} \times \{1\}\) is an orientation of this graph. One says that \((\widetilde{\mathrm{C}}, \mathrm{A})\) is the oriented graph associated with the quiver C.

If \(f\) is an arrow of C, one also says that \((f,1)\) is the oriented edge of \(\widetilde{C}\) associated with \(f\) and that the pair \(\{(f,1),(f,-1)\}\) is the edge of \(\widetilde{C}\) associated with \(f\) .

If C' is a subquiver of C, the directed graph \(\widetilde{C}'\) is a directed subgraph of \(\widetilde{C}\) .

There exists a unique quiver morphism \(\varphi\) from C to the quiver underlying \(\widetilde{\mathrm{C}}\) such that \(\varphi(f) = (f,1)\) for every arrow \(f\) of C; it is an isomorphism from C onto the quiver associated with the oriented graph ( \(\widetilde{\mathrm{C}}\) , A).

